\subsection{“视界零知识—幺正切片”范式：可交换极分解、联合谱梳与临界线集中}\label{subsec:op-algebra-hzom-unitary-slice-rh}
本小节在小节 \ref{subsec:op-algebra-modular-zk-index-fkdet} 的设定下，
把正文 HZOM（定义 \ref{def:xi-hzom}）所用的“视界算子族 \(\Rightarrow\) 行列式 \(\Rightarrow\) 零点几何”模板，
提升为一个更强的算子代数范式：在不可见子空间上若存在与可见粗粒化相容的\emph{可交换极分解}，
则临界线集中与零点局部统计可被完全转写为 \((H,U)\) 的联合谱数据；
并且 Li 判据可在该范式中被算子化为完全正迹泛函所诱导的离散凸性公理。

\paragraph{设置：不可见子空间与行列式零点事件}
令 \((\mathcal{M},\tau)\) 为带忠实正规迹的有限 von Neumann 代数，\(\mathcal{N}\subset\mathcal{M}\) 为“视界可见”子代数，
并设存在迹保持条件期望 \(E:\mathcal{M}\to\mathcal{N}\)。
记 GNS 空间 \(L^2(\mathcal{M},\tau)\) 与子空间 \(L^2(\mathcal{N},\tau)\)，并令
\[
\mathfrak{H}_{\mathrm{inv}}:=L^2(\mathcal{M},\tau)\ominus L^2(\mathcal{N},\tau)
\]
为不可见子空间。
设 \(\{K_s\}_{\Re(s)>\tfrac12}\subset\mathcal{M}\) 为一族随 \(s\) 解析变化的元，并定义（Fuglede--Kadison）视界行列式
\[
D(s):=\Delta_\tau(I-K_s)\ \in\ [0,\infty).
\]
由命题 \ref{prop:op-algebra-fkdet-properties}，有
\[
D(s)=0\quad\Longleftrightarrow\quad I-K_s\ \text{在 }\mathcal{M}\text{ 中不可逆}
\quad\Longleftrightarrow\quad 1\in \mathrm{Spec}_{\mathcal{M}}(K_s).
\]
当 \(K_s\) 在某 Hilbert 空间表示下为迹类算子时，也可用 Fredholm 行列式 \(\det(I-K_s)\) 作为复解析对象；
其零点事件与 \(I-K_s\) 的不可逆性一致，从而与 \(D(s)\) 的消失事件一致。
本小节关于“零点位置/重数/间距”的叙述均以该不可逆性事件为核心语义。

\begin{theorem}[“视界零知识—幺正切片”充要范式：可交换极分解迫使临界线集中]\label{thm:op-algebra-hzom-commuting-polar-forces-critical-line}
\leanverified{paper\_op\_algebra\_hzom\_commuting\_polar\_forces\_critical\_line}
在上述设定下，进一步假设在不可见子空间 \(\mathfrak{H}_{\mathrm{inv}}\) 上存在正自伴算子 \(H\ge 0\) 与幺正算子 \(U\)，满足
\[
[H,U]=0,
\qquad
K_s\big|_{\mathfrak{H}_{\mathrm{inv}}}
=\exp\!\bigl(-(s-\tfrac12)H\bigr)\,U\qquad(\Re(s)>\tfrac12),
\]
并且 \(H\) 在 \(\mathfrak{H}_{\mathrm{inv}}\) 上具有严格下界（等价地：对任意 \(\sigma>\tfrac12\)，\(\|\exp(-(\sigma-\tfrac12)H)\|<1\)）。
则 \(D(s)\) 在右半平面 \(\Re(s)>\tfrac12\) 内无零点。
若再假定反射对称
\[
\boxed{\ D(s)=\overline{D(1-\overline{s})}\ },
\]
则 \(D(s)\) 的全部“非平凡零点事件”均被压缩到临界线 \(\Re(s)=\tfrac12\)。
\end{theorem}
\begin{proof}
取任意 \(\sigma>\tfrac12\)。由假设，\(\exp(-(\sigma-\tfrac12)H)\) 在 \(\mathfrak{H}_{\mathrm{inv}}\) 上为严格收缩，而 \(U\) 为幺正，且二者可交换。
因此
\[
\left\|K_{\sigma+it}\big|_{\mathfrak{H}_{\mathrm{inv}}}\right\|
=\left\|\exp\!\bigl(-(\sigma-\tfrac12)H\bigr)\right\|
<1.
\]
从而 \(1\notin\mathrm{Spec}(K_{\sigma+it}|_{\mathfrak{H}_{\mathrm{inv}}})\)，并由不可逆性语义得到 \(I-K_{\sigma+it}\) 在不可见部分可逆；
由于可见部分对零知识接口仅提供可逆压缩（见小节 \ref{subsec:op-algebra-modular-zk-index-fkdet} 的“可见闭包”设定），得 \(I-K_{\sigma+it}\) 在 \(\mathcal{M}\) 中可逆，故 \(D(\sigma+it)\neq 0\)。
由反射对称再得左半平面 \(\Re(s)<\tfrac12\) 的零自由，从而非平凡零点事件只能发生在临界线。证毕。
\end{proof}

\begin{corollary}[临界零点的显式谱参数化：零点高度由 \texorpdfstring{$(H,U)$}{(H,U)} 的联合谱全决定]\label{cor:op-algebra-hzom-critical-zeros-joint-spectrum-comb}
\leanverified{paper\_op\_algebra\_hzom\_critical\_zeros\_joint\_spectrum\_comb}
在定理 \ref{thm:op-algebra-hzom-commuting-polar-forces-critical-line} 的设定下，
设 \(H\) 与 \(U\) 在 \(\mathfrak{H}_{\mathrm{inv}}\) 上强可交换，从而存在联合谱测度 \(\mathsf E(\dd\lambda,\dd\theta)\)（其中 \(\lambda\ge 0\)、\(\theta\in(-\pi,\pi]\)）使
\[
H=\int \lambda\,\mathsf E(\dd\lambda,\dd\theta),
\qquad
U=\int e^{i\theta}\,\mathsf E(\dd\lambda,\dd\theta).
\]
则在临界切片 \(s=\tfrac12+it\) 上，
\[
D(\tfrac12+it)=0
\quad\Longleftrightarrow\quad
\exists\,(\lambda,\theta)\ \text{在联合谱支撑上使}\ 1-e^{-it\lambda}e^{i\theta}=0,
\]
从而零点高度具有显式参数化
\[
\boxed{\ t=\frac{\theta+2\pi k}{\lambda}\qquad(\lambda>0,\ k\in\mathbb{Z}).\ }
\]
特别地：
\begin{enumerate}
  \item \textbf{（纤维重数）}在有限维/离散谱情形，零点的代数重数等于满足同一 \((\lambda,\theta)\) 的联合谱投影之迹维数；
  \item \textbf{（联合谱梳）}对固定尺度 \(\lambda\) 与相位 \(\theta\)，对应的零点列为等间距谱梳，其相邻间距严格等于 \(2\pi/\lambda\)。
\end{enumerate}
\end{corollary}
\begin{proof}
当 \(s=\tfrac12+it\) 时，\(K_s|_{\mathfrak{H}_{\mathrm{inv}}}=e^{-itH}U\) 为幺正算子。
由强可交换性，联合谱定理给出
\(
\mathrm{Spec}(e^{-itH}U)=\{e^{-it\lambda}e^{i\theta}:(\lambda,\theta)\in\mathrm{supp}(\mathsf E)\}
\)。
于是 \(1\in\mathrm{Spec}(e^{-itH}U)\) 当且仅当存在联合谱点使 \(e^{-it\lambda}e^{i\theta}=1\)，解得所述参数化。
有限维情形的重数结论由特征值 \(1\) 的谱投影维数（或其归一化迹）给出；间距结论对固定 \((\lambda,\theta)\) 直接由参数化得到。证毕。
\end{proof}

\begin{corollary}[临界零点简单性的联合谱判据：由视界相位的纤维简并完全决定]\label{cor:op-algebra-hzom-simple-zeros-joint-spectrum}
\leanverified{paper\_op\_algebra\_hzom\_simple\_zeros\_joint\_spectrum}
在推论 \ref{cor:op-algebra-hzom-critical-zeros-joint-spectrum-comb} 的离散谱口径下，
临界线上零点 \(s=\tfrac12+it\) 为简单零点当且仅当满足 \(e^{-it\lambda}e^{i\theta}=1\) 的联合谱投影在 \((\mathcal{M},\tau)\) 的迹意义下为一维（无纤维简并）。
等价地：对几乎处处 \(\lambda\)，幺正相位 \(U\) 在 \(H=\lambda\) 的谱纤维上具有简单谱。
\end{corollary}

\begin{corollary}[零点间距的尺度下界：固定尺度梳的统一控制]\label{cor:op-algebra-hzom-zero-spacing-scale-lower-bound}
\leanverified{paper\_op\_algebra\_hzom\_zero\_spacing\_scale\_lower\_bound}
在推论 \ref{cor:op-algebra-hzom-critical-zeros-joint-spectrum-comb} 的设定下，
对任意固定联合谱点 \((\lambda,\theta)\)（\(\lambda>0\)）所生成的谱梳零点列，其相邻间距恒为 \(2\pi/\lambda\)。
因此若 \(H\) 的联合谱在迹意义下本质上支撑于 \(\lambda\le \lambda_{\max}<\infty\)，则每一条固定尺度谱梳均满足统一下界
\[
\boxed{\ \Delta t_{\min}(\lambda,\theta)=\frac{2\pi}{\lambda}\ \ge\ \frac{2\pi}{\lambda_{\max}}.\ }
\]
从而在该范式中，若要在\emph{同一尺度分支}上产生任意小的零点间距，必须要求不可见生成元 \(H\) 具有无界尺度谱（\(\lambda_{\max}=\infty\)）。
\end{corollary}

\begin{proposition}[Li 判据的算子化：Li 系数等价于完全正迹泛函的离散凸性]\label{prop:op-algebra-li-criterion-operatorized}
\leanverified{paper\_op\_algebra\_li\_criterion\_operatorized}
在定理 \ref{thm:op-algebra-hzom-commuting-polar-forces-critical-line} 的设定下，
进一步假设 \(K_s\) 在某个可迹化口径下可定义 Fredholm 行列式 \(\det(I-K_s)\)，并取处处非零的解析正规化因子 \(G(s)\) 使完成化对象
\[
\Xi(s):=G(s)\,\det(I-K_s)
\]
满足 \(s\leftrightarrow 1-s\) 的反射对称。
令 \(\{\Lambda_n\}_{n\ge 1}\) 为 \(\Xi\) 的 Li 系数（按 Li 判据的标准生成定义）。
则存在一个完全正迹泛函 \(\mathfrak{L}:\mathcal{M}\to\RR\)（可取为 \(\mathfrak{L}=\tau\)）与一族仅依赖 \(n\) 的非交换多项式 \(P_n\)（作用于 \(e^{-H}\) 与 \(U\)，约定 \(P_0:=0\)），使得对所有 \(n\ge 1\) 有
\[
\boxed{\;
\Lambda_n
=\mathfrak{L}\!\bigl(P_{n+1}(e^{-H},U)\bigr)
-2\,\mathfrak{L}\!\bigl(P_{n}(e^{-H},U)\bigr)
+\mathfrak{L}\!\bigl(P_{n-1}(e^{-H},U)\bigr).\;}
\]
因此，Li 判据的非负性条件 \(\Lambda_n\ge 0\)（\(\forall n\)）等价于数列
\(
a_n:=\mathfrak{L}(P_n(e^{-H},U))
\)
在离散意义下的凸性：其二阶差分处处非负。
\end{proposition}
\begin{proof}[证明要点]
Li 系数由 \(\log \Xi\) 的局部生成展开给出。
在 \(\det(I-K_s)\) 可用时，有
\(
\log\det(I-K_s)=-\sum_{m\ge 1}\frac{1}{m}\Tr(K_s^m)
\)
（在收敛域内），而在 \([H,U]=0\) 下 \(K_s^m=e^{-m(s-\tfrac12)H}U^m\)。
将该展开与完成化因子 \(\log G(s)\) 合并，并把生成式逐项重排，即得到 Li 系数可表为有限个 \(\Tr(e^{-aH}U^b)\) 的线性组合；
这些量等价于对 \((e^{-H},U)\) 的非交换多项式取迹。
把这些线性组合按 \(n\) 重新积分为一列“原函数”值 \(a_n\)（即令 \(a_n\) 的二阶差分等于 \(\Lambda_n\)），并用非交换多项式 \(P_n\) 吸收该积分所需的组合系数，即得所述表示。
由于 \(\mathfrak{L}\) 完全正（例如 \(\tau\)），故 Li 判据的非负性等价于 \(a_n\) 的离散凸性。证毕。
\end{proof}

\begin{proposition}[共同幺正扩张判据：统一扩张存在性推出可交换极分解]\label{prop:op-algebra-common-unitary-dilation-criterion}
\leanverified{paper\_op\_algebra\_common\_unitary\_dilation\_criterion}
设 \(\{K_s\}_{\Re(s)>\tfrac12}\) 为视界核族，并假定对每个 \(\sigma>\tfrac12\) 与几乎处处 \(t\)，
\(
K_{\sigma+it}\big|_{\mathfrak{H}_{\mathrm{inv}}}
\)
为严格收缩。
若存在一个共同扩张空间 \(\widetilde{\mathfrak H}\supset\mathfrak{H}_{\mathrm{inv}}\) 与一族幺正算子 \(\{\widetilde U_{\sigma,t}\}\subset\mathcal{B}(\widetilde{\mathfrak H})\)，使得对所有 \((\sigma,t)\) 都有
\[
K_{\sigma+it}\big|_{\mathfrak{H}_{\mathrm{inv}}}
=P_{\mathrm{inv}}\,\widetilde U_{\sigma,t}\big|_{\mathfrak{H}_{\mathrm{inv}}},
\]
并且该扩张在 \(s\leftrightarrow 1-s\) 的对偶共轭对称下闭合（以模共轭实现），
则在同构意义下可推出存在可交换极分解表示
\(
K_s\big|_{\mathfrak{H}_{\mathrm{inv}}}=e^{-(s-\tfrac12)H}U
\)
（其中 \([H,U]=0\)），从而定理 \ref{thm:op-algebra-hzom-commuting-polar-forces-critical-line} 的临界线集中结论成立。
反之，若存在稳定的离临界线零点事件，则上述“共同幺正扩张与对偶闭合”不可同时成立。
\end{proposition}
\begin{proof}[证明要点]
每个固定 \((\sigma,t)\) 的严格收缩均可由 Sz.-Nagy 型定理得到幺正扩张；
“共同扩张”的假设强制这些扩张在同一谱模型下兼容，并在对偶闭合下携带 \(s\leftrightarrow 1-s\) 的反射结构。
在该兼容性下，可把 \((\sigma,t)\)-依赖吸收为一个正生成元 \(H\) 的指数衰减与一个边界幺正 \(U\) 的相位演化；
对偶闭合进一步强制 \(H\) 与 \(U\) 在不可见部分强可交换，从而得到所述极分解。
若存在离临界线零点，则右/左半平面出现不可逆性事件，与“统一扩张 + 对偶闭合 \(\Rightarrow\) 可逆压缩”的结构矛盾。证毕。
\end{proof}

\begin{proposition}[对数 Sobolev 常数控制零点束缚带：从混合不等式推出临界线局域化]\label{prop:op-algebra-logsobolev-zero-band}
\leanverified{paper\_op\_algebra\_logsobolev\_zero\_band}
设存在作用于 \(\mathfrak{H}_{\mathrm{inv}}\) 的完全正、迹保持半群 \((\mathcal{T}_u)_{u\ge 0}\) 与幺正扰动 \(U\)，使得
\[
K_s\big|_{\mathfrak{H}_{\mathrm{inv}}}
=\int_{0}^{\infty}e^{-(s-\tfrac12)u}\,\mathcal{T}_u(U)\,\dd u,
\]
且 \((\mathcal{T}_u)\) 满足维数无关的对数 Sobolev 不等式，其常数记为 \(\alpha>0\)（采用正文中 \(\alpha\) 越大混合越强的约定；见推论 \ref{cor:pom-fold-factor-chain-derived-invariants} 的超收缩表述）。
则存在显式常数 \(c>0\)（仅依赖正规化）使得视界行列式在带外区域
\[
|\Re(s)-\tfrac12|\ >\ \frac{c}{\alpha}
\]
无零点事件。
特别地，若在尺度极限中 \(\alpha\) 不退化（保持 \(\inf \alpha>0\)），则零点事件被强制压缩至临界线邻域；
反之，任何稳定存在的离线零点事件都会迫使相应尺度上的对数 Sobolev 常数崩塌（\(\alpha\downarrow 0\)）。
\end{proposition}
\begin{proof}[证明要点]
对数 Sobolev 不等式与 Gross 超收缩等价：存在 \(q(u)=1+e^{2\alpha u}\) 使 \(\|\mathcal{T}_u f\|_{q(u)}\le \|f\|_2\)。
在 \(U\) 位于不可见部分且满足标准中心化/正交条件时，超收缩可提升为对 \(\|\mathcal{T}_u(U)\|\) 的指数衰减界；
与 Laplace 权 \(e^{-(\Re(s)-1/2)u}\) 叠乘后，得到 \(K_s\) 的严格收缩性，从而 \(1\notin\mathrm{Spec}(K_s)\) 并排除零点事件。
常数 \(c\) 由上述两类指数衰减的拼接给出；对偶反射再把右半平面的零自由传递到左半平面。证毕。
\end{proof}

\endinput
